\chapter{Diseño de la secuencia didáctica 5E}\label{cap:propuesta}

\textit{Este capítulo desarrolla la Fase 2 de la metodología (Capítulo 1): a partir de la síntesis de la revisión (Capítulo 2) se diseña la secuencia didáctica 5E, decidiendo ---con base en sus resultados--- qué contenidos incluir, qué estrategias emplear, qué ideas erróneas abordar y cómo evaluar la comprensión. La estructura, los objetivos, la secuenciación y el encaje curricular ya están definidos; el \textbf{desarrollo detallado} de cada actividad ---con su contenido, sus materiales y su evaluación--- se aborda en el Capítulo 4 (§4.2--§4.6).}

\section{Tema, método y delimitación de la propuesta}\label{sec:tema}

\textbf{Tema.} La secuencia \textbf{parte de un problema} planteado desde la primera sesión ---\textbf{¿podemos determinar la edad y el tamaño del universo? ¿Cómo lo sabe la ciencia?}--- y avanza hasta resolverlo, culminando en esa \textbf{determinación de la edad y el tamaño del universo}. El eje no es solo el resultado, sino \textbf{el proceso, la metodología y las herramientas} que conducen a esa medida: de los espectros y el corrimiento al rojo a la ley de Hubble-Lemaître, de ésta a la constante de Hubble ($H_{0}$) y, finalmente, a la \textbf{edad del universo} ($t \approx 1/H_{0}$) y al \textbf{tamaño del universo observable}. La expansión acelerada y la energía oscura se introducen \textbf{solo como matiz} en la fase Elaborate ---la razón por la que la estimación $t \approx 1/H_{0}$ es una primera aproximación---, no como núcleo. La materia oscura y la radiación cósmica de fondo se emplean como contenidos de apoyo. Queda \textbf{fuera del alcance} el formalismo de la relatividad general y las funciones de distancia cosmológicas.

\textbf{Eje conceptual: el espacio y el tiempo, de lo local a lo cosmológico.} Bajo la determinación de la edad y el tamaño late un objetivo conceptual más profundo: \textbf{revisar las nociones de espacio y de tiempo} al pasar de la escala cotidiana a la astronómica y cosmológica. La secuencia hace explícito ese tránsito y, en cada fase, esos conceptos se \textbf{redefinen, matizan, entienden o calculan}:

\begin{itemize}
\item \textbf{Espacio.} De la distancia local que se mide con una regla (aula, planeta) a las \textbf{distancias astronómicas} (el año-luz, una unidad que ya entrelaza espacio y tiempo) y, finalmente, a la idea contraintuitiva de un \textbf{espacio que se expande}: no son las galaxias las que se mueven por el espacio, sino el propio espacio el que se estira.
\item \textbf{Tiempo.} Del tiempo cotidiano a las \textbf{escalas de miles de millones de años}, y a una consecuencia de la velocidad finita de la luz: \textbf{mirar lejos es mirar al pasado} (\textit{look-back time}). Determinar la edad del universo es, literalmente, medir un tiempo.
\item \textbf{La unión de ambos.} A escala cosmológica, espacio y tiempo dejan de ser independientes: la distancia se infiere del corrimiento al rojo y la edad se obtiene de una velocidad y una distancia ($t \approx 1/H_{0}$). El alumnado descubre que \textbf{la misma velocidad de la luz que limita lo que podemos ver (tamaño) fija hasta cuándo podemos mirar hacia atrás (edad)}.
\end{itemize}

Así, cada fase no solo calcula un número: \textbf{reconfigura las intuiciones del alumnado sobre el espacio y el tiempo} ---qué se puede medir, con qué unidades y por qué lo cotidiano no basta a escala cósmica---.

\textbf{Precisión conceptual: qué entendemos por «edad», «tamaño» y «universo».} Conviene delimitar estos tres términos, porque son fuente habitual de confusión y la secuencia podría introducir ideas erróneas si no se explicitan:

\begin{itemize}
\item \textbf{Universo.} Se distingue el \textbf{universo observable} ---la región de la que nos puede llegar información, limitada por la edad del cosmos y la velocidad de la luz--- del \textbf{universo en su conjunto}, del que no conocemos ni borde ni centro y que podría ser infinito. En el aula \textbf{siempre estimamos el observable}; debe evitarse la imagen de un universo con un «tamaño total» definido y un borde.
\item \textbf{Edad.} Es el \textbf{tiempo transcurrido desde el Big Bang} ($\sim$13\,800 millones de años). El Big Bang no es una explosión \textit{en} el espacio, sino el origen del propio espacio-tiempo. La estimación escolar $t \approx 1/H_{0}$ (tiempo de Hubble, $\sim$14\,000 millones de años por \textbf{extrapolación lineal} de la tasa de expansión actual) es una \textbf{primera aproximación} ---de ahí el matiz de la fase Elaborate: la edad precisa depende de toda la historia de la expansión---.
\item \textbf{Tamaño.} Sutileza que conviene hacer explícita para \textbf{no enseñar sin querer un error}: el radio del universo observable \textbf{no es} $c \times$ edad ($\approx$ 13\,800 millones de años-luz), sino \textbf{bastante mayor} ($\sim$46\,000 millones de años-luz), porque el espacio se ha estirado mientras la luz viajaba. En bachillerato se estima el \textbf{orden de magnitud} ---el radio de Hubble, $c/H_{0}$--- dejando claro que el cálculo exacto del radio observable exige integrar la expansión (fuera del alcance). Y, de nuevo, «tamaño del observable» $\neq$ «tamaño del universo total» (desconocido).
\end{itemize}

En suma, la secuencia enseña a \textbf{estimar el universo observable} y a tratar «edad» y «tamaño» como magnitudes \textbf{inferidas y aproximadas}, no como datos absolutos, lo que refuerza el objetivo transversal sobre la naturaleza de la ciencia (OA7).

\textbf{Método.} Se adopta el \textbf{modelo 5E} \parencite{bybee2006} como marco de diseño instruccional. Su elección se fundamenta en la revisión: la secuenciación observacional basada en datos es la vía más honesta científicamente para construir la determinación de la edad y el tamaño a partir de las observaciones ---e introducir, ya como matiz, la energía oscura--- (§2.4.2), y el 5E la estructura de forma natural (indagación $\rightarrow$ formalización $\rightarrow$ extensión $\rightarrow$ evaluación). Cada fase integra elementos concretos extraídos de la revisión (mapeo de §2.4.4) y materializa las cinco decisiones de diseño (§2.6).

\textbf{Contexto.} La propuesta se diseña \textbf{para 2º de Bachillerato en el marco de Pearson Edexcel International Advanced Level (IAL) Physics}, concretamente la \textbf{Unit~5.6 ``Astrophysics and Cosmology''} (en el libro de texto, subtema \textit{11B ``Space''} del \textit{Topic~11}), donde estos contenidos son \textbf{currículo evaluable} (§2.5). Es el contexto en el que el autor imparte docencia y en el que la propuesta encaja con mayor precisión, lo que facilita su implementación en el curso 2026-2027. El detalle de destinatarios y prerrequisitos se recoge en §3.3.

\section{Fundamentos del modelo 5E y comparación con el enfoque tradicional}\label{sec:fundamentos5e}

\textbf{Fundamentos.} El modelo 5E fue formalizado por Rodger W. Bybee y el \textit{Biological Sciences Curriculum Study} (BSCS) \parencite{bybee2006} como evolución del \textit{learning cycle} de \textcite{atkin1962}. Se apoya en una \textbf{concepción constructivista del aprendizaje} ---el conocimiento se construye de forma activa sobre las ideas previas \parencite{piaget1964} y en interacción social \parencite{vygotsky1978}--- y en la \textbf{teoría del cambio conceptual} \parencite{posner1982}, según la cual aprender ciencia exige que el alumnado confronte y reestructure sus concepciones alternativas, no que reciba información ya elaborada. Esto es especialmente pertinente aquí, porque la cosmología está plagada de ideas erróneas resistentes (§2.3.4).

El modelo organiza la enseñanza en \textbf{cinco fases}:

\begin{itemize}
\item \textbf{Engage (implicar):} despertar el interés y aflorar las ideas previas mediante una situación o pregunta motriz.
\item \textbf{Explore (explorar):} el alumnado indaga con materiales o datos ---sin recibir aún la explicación formal--- y genera experiencias comunes sobre las que construir.
\item \textbf{Explain (explicar):} se formaliza el concepto; el alumnado articula lo explorado y el docente introduce la terminología y los modelos.
\item \textbf{Elaborate (elaborar):} se aplica y transfiere lo aprendido a situaciones nuevas, profundizando o matizando.
\item \textbf{Evaluate (evaluar):} se valora la comprensión ---de forma continua y final--- frente a los objetivos.
\end{itemize}

En este tránsito, el papel del docente se desplaza del de transmisor al de \textbf{guía de la indagación}, y el del alumnado del de receptor al de \textbf{agente que mide, infiere y explica}.

\textbf{Eficacia y encaje.} La investigación asocia el 5E con mejoras en la comprensión conceptual y en la motivación frente a la enseñanza expositiva \parencite{bybee2006}. Su lógica ---primero la experiencia y los datos, después el modelo--- coincide con la \textbf{vía observacional} que la revisión identificó como la más honesta científicamente para esta materia (§2.4.2) y con la necesidad de \textbf{confrontar de forma explícita las ideas erróneas} (§2.4.3). Por todo ello se adopta como marco de la secuencia.

\textbf{Comparación con el enfoque tradicional (del libro de texto).} El desarrollo habitual de T11·11B en el libro de Edexcel es \textbf{expositivo}: el docente presenta los conceptos (distancias, Doppler, corrimiento al rojo, ley de Hubble-Lemaître, edad y destino del universo) y el alumnado los aplica en ejercicios de tipo examen. La secuencia 5E \textbf{reordena el mismo contenido evaluable} en torno a la indagación. La comparación ---con sus ventajas y sus contrapartidas--- se resume en la Tabla~\ref{tab:comparacion}.

\begin{table}[ht]
\centering
\caption{Comparación entre el enfoque tradicional y la secuencia 5E propuesta.}\label{tab:comparacion}
\begin{tabularx}{\textwidth}{|>{\raggedright\arraybackslash}p{2.7cm}|>{\raggedright\arraybackslash}X|>{\raggedright\arraybackslash}X|}
\hline
\rowcolor{naranja}
{\color{white}\textbf{Dimensión}} & {\color{white}\textbf{Enfoque tradicional (libro de texto)}} & {\color{white}\textbf{Secuencia 5E propuesta}} \\ \hline
Punto de partida & Exposición directa del contenido & Pregunta motriz e ideas previas (Engage) \\ \hline
Papel del alumnado & Receptor; aplica fórmulas dadas & Activo: mide, construye el diagrama de Hubble, infiere $H_{0}$ y la edad \\ \hline
Papel del docente & Transmisor de contenido & Guía de la indagación \\ \hline
Ley de Hubble-Lemaître / $H_{0}$ & Se da la fórmula $v = H_{0}d$ y se practica & Emerge de los datos del propio alumnado \\ \hline
Edad del universo & Dato que se memoriza ($\sim$13,8 Ga) & Resultado que el alumnado calcula, interpreta y matiza \\ \hline
Ideas erróneas & Quedan latentes & Se afloran y confrontan de forma explícita (§3.5) \\ \hline
Espacio y tiempo & Definiciones dadas & Se reconstruyen de lo local a lo cósmico (§3.1) \\ \hline
Naturaleza de la ciencia & Hechos cerrados & Proceso de inferencia visible (OA7) \\ \hline
Motivación & Variable & Alta (relevancia y protagonismo del alumnado) \\ \hline
Evaluación & Examen de contenidos & Comprensión conceptual y del proceso; medida pre/post \\ \hline
Coste & Menor tiempo y preparación & Mayor tiempo de aula y preparación docente \\ \hline
\end{tabularx}
\end{table}

Esta comparación es coherente con la revisión: la enseñanza expositiva, aunque \textbf{eficiente para cubrir temario}, produce ganancias conceptuales limitadas y deja intactas las ideas erróneas (§2.4.3), mientras que la indagación basada en datos genera \textbf{más comprensión} a costa de más tiempo y recursos (§2.4.2; barreras en §2.3.4). La secuencia 5E asume conscientemente ese intercambio ---algo de eficiencia a cambio de profundidad y cambio conceptual---, justificado porque estos contenidos son precisamente los \textbf{más propensos a ideas erróneas} y los que peor resisten un tratamiento superficial.

\section{Contexto y destinatarios}\label{sec:destinatarios}

\begin{itemize}
\item \textbf{Nivel:} 2º de Bachillerato.
\item \textbf{Marco curricular:} Pearson Edexcel IAL Physics, \textbf{Unit~5.6 --- Astrophysics and Cosmology} (\textit{Topic~11} del libro, subtema 11B ``Space''). Los objetivos de la secuencia se alinean con sus estándares evaluables (redshift, $z = \Delta\lambda/\lambda \approx v/c$, $v = H_{0}d$, edad y destino del universo, materia y energía oscura). Ver §2.5.
\item \textbf{Prerrequisitos del alumnado} (cubiertos en Edexcel IAL Año 1 y en Topic 11): concepto de onda y efecto Doppler, conservación de la energía, campo gravitatorio y órbitas, ajuste de una recta.
\item \textbf{Recursos disponibles:} los \textbf{habituales del aula}; el \textbf{recurso didáctico central es la web interactiva}, de uso \textbf{offline} (sin instalación ni conexión), por lo que basta con \textbf{un dispositivo con navegador por pareja} (tablet u ordenador) y un \textbf{proyector} para las puestas en común. Cada actividad interactiva dispone de una \textbf{alternativa en papel} equivalente (equidad), lo que hace la secuencia reproducible en cualquier aula.
\end{itemize}

\section{Objetivos de aprendizaje}\label{sec:objetivos5e}

\textbf{Objetivo general.} Que el alumnado \textbf{comprenda y aplique el proceso por el que la ciencia determina la edad y el tamaño del universo} ---de la medida de distancias y del corrimiento al rojo a la ley de Hubble-Lemaître y a la estimación de $t \approx 1/H_{0}$---, revisando en ese recorrido sus nociones de espacio y de tiempo.

Los \textbf{objetivos de aprendizaje} (OA) se alinean con los estándares evaluables de Edexcel IAL Physics (Unit~5.6; subtema 11B del \textit{Topic~11}), se ubican en la fase 5E donde se trabajan y señalan la idea errónea que abordan (Tabla~\ref{tab:oa}).

\setlength{\tabcolsep}{3pt}
\begin{longtable}{>{\raggedright\arraybackslash}p{0.9cm} >{\raggedright\arraybackslash}p{5.3cm} >{\raggedright\arraybackslash}p{3.0cm} >{\raggedright\arraybackslash}p{1.5cm} >{\raggedright\arraybackslash}p{3.6cm}}
\caption{Objetivos de aprendizaje, encaje curricular, fase 5E e idea errónea que abordan.}\label{tab:oa}\\
\rowcolor{naranja}
{\color{white}\textbf{Código}} & {\color{white}\textbf{Objetivo de aprendizaje}} & {\color{white}\textbf{Contenido Edexcel 5.6 (11B)}} & {\color{white}\textbf{Fase 5E}} & {\color{white}\textbf{Idea errónea que aborda}} \\ \hline
\endfirsthead
\rowcolor{naranja}
{\color{white}\textbf{Código}} & {\color{white}\textbf{Objetivo de aprendizaje}} & {\color{white}\textbf{Contenido Edexcel 5.6 (11B)}} & {\color{white}\textbf{Fase 5E}} & {\color{white}\textbf{Idea errónea que aborda}} \\ \hline
\endhead
OA1 & Determinar distancias astronómicas mediante paralaje trigonométrico y candelas estándar, y reconocer el año-luz como unidad que entrelaza espacio y tiempo. & Distancias: paralaje y candelas estándar & Explore & Dificultad con las escalas espaciales \\ \hline
OA2 & Interpretar el corrimiento al rojo de los espectros galácticos como estiramiento de la longitud de onda por la expansión del espacio, distinguiéndolo del efecto Doppler ordinario. & Efecto Doppler y corrimiento al rojo ($z = \Delta\lambda/\lambda \approx v/c$) & Explore $\rightarrow$ Explain & Todo corrimiento al rojo es efecto Doppler \\ \hline
OA3 & Construir e interpretar el diagrama de Hubble y determinar la constante de Hubble $H_{0}$ a partir de datos (reales o simulados). & Ley de Hubble-Lemaître ($v = H_{0}d$) & Explore $\rightarrow$ Explain & Expansión = movimiento propio de las galaxias \\ \hline
OA4 & Estimar la edad del universo ($t \approx 1/H_{0}$) y el tamaño del universo observable, e interpretar el significado físico de la medida. & Edad del universo; tamaño observable (nivel de libro de texto) & Explain $\rightarrow$ Elaborate & Big Bang como explosión puntual; universo con centro/borde \\ \hline
OA5 & Reconocer que $t \approx 1/H_{0}$ es una extrapolación lineal simplificada de la tasa actual y que la edad precisa depende de la historia de la expansión (que varía: la gravedad frena, la energía oscura acelera), lo que matiza la medida. & Destino del universo y constante de Hubble; energía oscura (nivel de libro de texto) & Elaborate & $t = 1/H_{0}$ como valor exacto; energía oscura como fuerza/sustancia \\ \hline
OA6 & Diferenciar la materia oscura de la energía oscura como conceptos distintos. & Materia oscura & Elaborate & Confusión materia oscura / energía oscura \\ \hline
OA7 & Valorar cómo la ciencia infiere magnitudes inaccesibles (edad, tamaño) a partir de observaciones, modelos y aproximaciones, reconociendo sus límites y que la ciencia \textbf{progresa mediante debates entre modelos y teorías} provisionales y sujetos a la evidencia. \textit{(Objetivo transversal --- naturaleza de la ciencia.)} & Transversal (5.6 en su conjunto) & Todas & Ciencia como conjunto de hechos cerrados \\ \hline
\end{longtable}

\section{Ideas erróneas que aborda la secuencia}\label{sec:ideaserroneas}

Se atacan de forma explícita las seis ideas erróneas nucleares (§2.3.4): Big Bang como explosión puntual; universo con centro o borde; expansión de los sistemas ligados; corrimiento al rojo como Doppler; confusión materia/energía oscura; energía oscura como «fuerza que empuja». \textit{(Cada fase indica cuál confronta.)}

\section{Temporalización}\label{sec:temporalizacion}

\textbf{Encaje en la programación real (curso 2026-2027).} La secuencia se ubica en \textbf{Y13 (equivalente a 2º de Bachillerato), tercer trimestre (abril de 2027)}, dentro de \textbf{T11 --- Astrophysics and Cosmology, subtema 11B «Space»} del \textit{scheme of work} de Pearson Edexcel IAL. Ocupa concretamente el hilo \textbf{distancias astronómicas $\rightarrow$ efecto Doppler y corrimiento al rojo $\rightarrow$ ley de Hubble-Lemaître $\rightarrow$ edad y destino del universo} (aprox. semanas 27--28 del calendario de Y13), una vez impartidos 11A (campos gravitatorios) y la parte estelar de 11B (radiación de cuerpo negro y diagrama H-R). No cubre, por tanto, todo el T11, sino el bloque de la \textbf{medida del universo}.

\textbf{Formato.} Clases de \textbf{45 min} (modelo de 6 clases/semana). La propuesta ocupa \textbf{siete sesiones}, compatible con la compresión del tercer trimestre (T11 dispone de $\sim$4 semanas y culmina con el examen de simulacro).

\begin{table}[ht]
\centering
\caption{Temporalización de la secuencia por sesiones y fase 5E.}\label{tab:temporalizacion}
\begin{tabularx}{\textwidth}{|>{\raggedright\arraybackslash}p{1.3cm}|>{\raggedright\arraybackslash}p{1.7cm}|>{\raggedright\arraybackslash}X|>{\raggedright\arraybackslash}p{3.2cm}|}
\hline
\rowcolor{naranja}
{\color{white}\textbf{Sesión}} & {\color{white}\textbf{Fase 5E}} & {\color{white}\textbf{Contenido / actividad}} & {\color{white}\textbf{Contenido Edexcel 11B}} \\ \hline
1 & Engage & El reto y los términos clave; aventura de la cosmología; paradoja de Olbers; cuestionario conceptual (pretest) & (motivación; enlaza con paralaje y candelas estándar) \\ \hline
2 & Explore & De la distancia (paralaje / candelas estándar) a la recesión (redshift): observar y medir & Distancias astronómicas; Doppler; redshift \\ \hline
3 & Explore & Construir el diagrama de Hubble y obtener $H_{0}$ (banda elástica + Tracker o datos reales) & $v = H_{0}d$ \\ \hline
4 & Explain & De $H_{0}$ a la edad ($t \approx 1/H_{0}$); la edad en perspectiva (calendario cósmico); corrimiento al rojo cosmológico vs. Doppler & Corrimiento al rojo; Hubble; edad del universo \\ \hline
5 & Elaborate & El \textbf{tamaño} del universo observable: radio de Hubble ($c/H_{0}$) frente al radio observable ($\sim$46\,000 Mal, de la historia de la expansión); derivación cualitativa con la CMB ($z \approx 1100$); transferencia a otros valores de $H_{0}$ & Tamaño del universo observable \\ \hline
6 & Elaborate & Los \textbf{límites y matices}: alcance de $1/H_{0}$; horizonte de sucesos; $\Lambda$CDM y tensión de Hubble; energía oscura (matiz) & Destino del universo; materia/energía oscura \\ \hline
7 & Evaluate & Cuestionario post; reconstrucción de la cadena y producto final & \textit{Formative} T11 \\ \hline
\end{tabularx}
\end{table}

\textit{Nota:} la parte estelar de 11B (cuerpo negro, Stefan-Boltzmann, Wien, diagrama H-R) se imparte aparte y no forma parte de la secuencia; sí se aprovechan \textbf{paralaje y candelas estándar} como herramientas de medida de \textbf{distancias} (el «espacio» del eje conceptual de §3.1).

\section{Estructura de la secuencia por fases}\label{sec:desarrollo}

\textbf{Pregunta motriz.} Toda la secuencia gira en torno a una pregunta que permanece visible en todas las sesiones: \textbf{«¿Podemos determinar la edad y el tamaño del universo? ¿Cómo lo sabe la ciencia?»}. Cada fase responde a una pregunta parcial que abre la siguiente, de modo que el conocimiento avanza como un \textbf{descubrimiento encadenado} ---de «esto no se puede saber» a «sí se puede, y así lo determinamos»--- y no como una suma de temas inconexos.

\begin{figure}[htbp]
\centering
\begin{tikzpicture}[
  font=\footnotesize,
  node distance=3.5mm,
  eq/.style={rectangle, rounded corners, draw=naranja, line width=0.6pt, fill=naranja!10, inner sep=1.4mm, align=center, text width=2.05cm, minimum height=1.3cm},
  ini/.style={rectangle, rounded corners, draw=naranja, line width=0.9pt, fill=naranja, text=white, inner sep=1.4mm, align=center, text width=1.7cm, minimum height=1.3cm},
  fin/.style={rectangle, rounded corners, draw=black!30, dashed, fill=black!5, inner sep=1.6mm, align=center, text width=6.0cm, minimum height=0.9cm},
  fl/.style={-{Stealth[length=1.9mm]}, draw=naranja, line width=0.85pt}
]
\node[ini] (q) {\textbf{¿Edad y tamaño?}};
\node[eq, right=of q] (z) {corrimiento al rojo\\[2pt]$z=\Delta\lambda/\lambda$};
\node[eq, right=of z] (v) {velocidad\\[2pt]$v=c\,z$};
\node[eq, right=of v] (h) {ley de Hubble-Lemaître\\[2pt]$v=H_{0}\,d$};
\node[eq, right=of h] (t) {edad\\[2pt]$t\approx 1/H_{0}$};
\node[eq, right=of t] (r) {tamaño\\[2pt]$R\approx c/H_{0}$};
\node[fin, below=5mm of h] (lim) {\textbf{límites:} energía oscura, tensión de Hubble y horizonte de sucesos};
\draw[fl](q)--(z);
\draw[fl](z)--(v);
\draw[fl](v)--(h);
\draw[fl](h)--(t);
\draw[fl](t)--(r);
\draw[fl, black!35, dashed](r.south) |- (lim.east);
\end{tikzpicture}
\caption{La \textbf{cadena de la medida}: de la observación (el corrimiento al rojo) a la edad y el tamaño del universo, y sus límites. Es el hilo conductor conceptual que recorre toda la secuencia.}\label{fig:cadena}
\end{figure}

\textbf{Hilo conductor: una progresión de conceptos.} La secuencia no acumula cifras, sino que se diseña como una \textbf{progresión de conceptos}: cada fase aborda la idea que la siguiente necesita, de modo que la cadena (Figura~\ref{fig:cadena}) ---del corrimiento al rojo a $H_{0}$, y de ahí a la edad y el tamaño--- \textbf{conduce a comprender el valor final}. En la web interactiva, esa progresión se muestra en cada página mediante un hilo que marca en qué punto del recorrido se encuentra el alumnado (§4.1). Entre fases, la \textbf{pregunta puente} actúa de bisagra reflexiva: al cerrar cada etapa se explicita \textit{qué sabíamos}, \textit{qué hemos averiguado} y \textit{qué queda abierto}, conectando el final de una sesión con el inicio de la siguiente. La Tabla~\ref{tab:correspondencia} hace explícita la correspondencia entre cada etapa, los conceptos que aborda y la forma de abordarlos.

\begin{table}[htbp]
\centering
\caption{Correspondencia entre cada etapa del ciclo 5E, los conceptos que aborda, la forma de abordarlos y su aportación a la comprensión del valor final (la edad y el tamaño del universo).}\label{tab:correspondencia}
\begin{tabularx}{\textwidth}{|>{\raggedright\arraybackslash}p{1.7cm}|>{\raggedright\arraybackslash}X|>{\raggedright\arraybackslash}X|>{\raggedright\arraybackslash}X|}
\hline
\rowcolor{naranja}
{\color{white}\textbf{Etapa}} & {\color{white}\textbf{Concepto(s) que aborda}} & {\color{white}\textbf{Cómo se aborda (5E)}} & {\color{white}\textbf{Qué aporta al valor final}} \\ \hline
Engage & Qué significan la edad y el tamaño; la enormidad del tiempo cósmico (como pregunta); la ciencia evoluciona con la tecnología. & Enmarcamos el reto, definimos los términos, recorremos la historia de la cosmología y provocamos con Olbers. & Motiva y plantea el problema: ¿se puede determinar? \\ \hline
Explore & Corrimiento al rojo, velocidad, distancia y el patrón entre velocidad y distancia. & Indagamos con datos: medimos $z$, construimos el diagrama de Hubble y registramos nuestra conjetura, sin interpretarla aún. & Obtiene la pendiente del diagrama ($H_{0}$), todavía sin interpretar. \\ \hline
Explain & La ley de Hubble-Lemaître ($v = H_{0}d$); qué es $H_{0}$ (el ritmo de expansión); el corrimiento al rojo como estiramiento del espacio; $1/H_{0}$; la escala del tiempo cósmico. & El alumnado explica primero; luego formalizamos la ley y razonamos: invertir $H_{0}$ es rebobinar la expansión; ponemos la edad en perspectiva (calendario cósmico). & Convierte $H_{0}$ en la \textbf{edad} ($t \approx 1/H_{0}$). \\ \hline
Elaborate & Tamaño observable; límites de la estimación; $\Lambda$CDM y tensión de Hubble. & Calculamos y contextualizamos, con el matiz de la energía oscura. & Obtiene el \textbf{tamaño} observable y reconoce sus límites. \\ \hline
Evaluate & El proceso completo de medida como objeto de comprensión. & Sintetizamos y lo explicamos a otra persona. & Consolida el entendimiento del valor final. \\ \hline
\end{tabularx}
\end{table}

\subsection{Diseño reproducible de la secuencia}\label{sec:diseno5e_reproducible}

\textbf{Un diseño reproducible.} El criterio fundamental del diseño no es distribuir contenidos en cinco apartados, sino lograr que el alumnado \textbf{explore y genere explicaciones preliminares antes de recibir la formalización científica}: por eso \textit{Explore} precede conceptualmente a \textit{Explain}, mientras que \textit{Elaborate} se reserva para \textbf{transferir} lo aprendido a situaciones nuevas y no para incorporar contenidos adicionales. Todas las actividades principales se integran en la plataforma web \textit{Cosmos5E} y pueden realizarse íntegramente en un navegador; los datos, gráficos y materiales esenciales se almacenan \textbf{localmente} junto con la aplicación, de modo que la implementación no dependa de servicios externos y sea \textbf{reproducible}, conservando las fuentes originales como referencias y elementos de trazabilidad. La distribución concreta en siete sesiones se detalla en la Temporalización (§\ref{sec:temporalizacion}).

\subsubsection{Criterios para la selección y adaptación de las actividades}

Las actividades no se seleccionaron únicamente por su atractivo visual o por su disponibilidad tecnológica. Se emplearon tres criterios de diseño:

\begin{enumerate}
    \item adecuación conceptual y matemática al nivel de 2.\textsuperscript{o} de Bachillerato;
    \item correspondencia con la función específica de la fase 5E en la que se incorpora la actividad;
    \item posibilidad de implementación digital, reproducible y basada, cuando resulte posible, en observaciones o datos astronómicos auténticos.
\end{enumerate}

Como antecedentes principales, se consideraron actividades educativas que reconstruyen la expansión cósmica a partir de datos observacionales. La actividad \textit{Measuring the Age of the Universe}, desarrollada por Las Cumbres Observatory \parencite{LCOAgeUniverse}, guía al alumnado desde la medición del corrimiento al rojo y de las distancias hasta la construcción del diagrama de Hubble, la estimación de $H_0$ y el cálculo aproximado de la edad del universo. Este recurso se tomó como referencia para la progresión observación--medida--representación--inferencia, pero se simplificó para evitar introducir simultáneamente la física de las supernovas Ia y el módulo de distancia, contenidos que no constituyen el objetivo central de la presente propuesta.

También se revisaron los materiales \textit{Cosmic Times} de NASA, en particular la actividad \textit{Determining the Nature, Size, and Age of the Universe} \parencite{NASACosmicTimes1929}, en la que el alumnado trabaja con observaciones y medidas de galaxias para construir una representación gráfica de la relación entre distancia y velocidad de recesión y, a partir de ella, discutir la expansión del universo. De este recurso se extrae la idea de presentar la expansión como un problema que debe inferirse a partir de evidencias observacionales antes de introducir su explicación formal. Asimismo, la propuesta de \textcite{cardona2016} constituye un antecedente directo para la reconstrucción escolar de la ley de Hubble a partir del análisis de datos astronómicos.

Estas actividades se adaptaron y no se reprodujeron literalmente para ajustarlas al hilo conceptual del presente trabajo. La finalidad no es que el alumnado repita un procedimiento previamente explicado, sino que descubra una regularidad en los datos y que posteriormente pueda interpretarla.

\subsection{Recorrido detallado por fases}\label{sec:recorrido-fases}

\subsubsection{Fase 1: Engage --- planteamiento del problema y elicitación de ideas previas}\label{sec:engage}

\paragraph{Pregunta inicial.} \emph{Tenemos luz procedente de galaxias situadas a millones de años-luz. ¿Podemos utilizar únicamente esa luz para averiguar si el universo cambia y estimar cuánto tiempo lleva evolucionando?}

La finalidad de esta fase es despertar el interés, identificar las concepciones iniciales del alumnado y establecer la necesidad de obtener evidencias. En consecuencia, durante \textit{Engage} no se explican todavía la ley de Hubble--Lemaître, la constante de Hubble ni el procedimiento empleado para estimar la edad del universo.

La sesión comienza con la aplicación individual de un cuestionario conceptual inicial (\textit{pretest}), cuyos resultados no se corrigen ni se muestran al alumnado en esta fase. El instrumento ---cuyos distractores recogen las ideas erróneas nucleares (§\ref{sec:ideaserroneas}) e incluye la opción «No lo sé / no lo tengo claro» para quien no lo tenga claro--- se describe en los instrumentos de evaluación y se aplica nuevamente, sin modificaciones sustantivas, durante \textit{Evaluate}, de modo que sirve de \textbf{línea base} para medir el cambio conceptual. Como situación provocadora se introduce después la \textbf{paradoja de Olbers} ---si el universo fuera eterno, estático y estuviera poblado uniformemente por estrellas, ¿por qué el cielo nocturno no sería completamente luminoso?---, planteada como interrogante y sin resolver todavía. El producto de la fase es el registro diagnóstico inicial del alumnado, que se contrasta explícitamente en \textit{Evaluate}.

\subsubsection{Fase 2: Explore --- indagación con observaciones y datos}\label{sec:explore}

La fase \textit{Explore} constituye el núcleo indagatorio de la secuencia y abarca dos sesiones. El alumnado trabaja primero con información espectral y, posteriormente, con datos de galaxias. La relación de Hubble--Lemaître no se presenta de forma previa, sino que debe emerger del análisis.

\paragraph{El corrimiento al rojo.} La web presenta una línea espectral cuya longitud de onda en reposo es conocida. El alumnado identifica su posición observada y calcula
\begin{equation}
z=\frac{\lambda_{\mathrm{obs}}-\lambda_0}{\lambda_0}.
\end{equation}
En esta etapa, $z$ se introduce únicamente como una medida del desplazamiento relativo de la línea. Para los objetos utilizados, todos en régimen de bajo corrimiento al rojo, se proporciona después la aproximación $v\simeq cz$, dejando explícito que su interpretación física y sus límites se discutirán en \textit{Explain}. Antes de cualquier interpretación, el alumnado responde: ¿hacia qué región del espectro se ha desplazado la línea?; ¿el desplazamiento es igual para todos los objetos?; ¿qué magnitud permite cuantificarlo?; ¿qué podría significar físicamente? La actividad de \textit{Cosmos5E}, basada en una línea H$\alpha$ interactiva, se mantiene como andamiaje inicial y se complementa con ejemplos de observaciones reales de procedencia documentada.

\paragraph{El diagrama de Hubble.} La segunda sesión utiliza un conjunto cerrado y versionado de galaxias cuyas distancias se han determinado mediante métodos independientes del corrimiento al rojo; el conjunto se incorpora a la plataforma en formato CSV/JSON para garantizar la reproducibilidad. El alumnado representa la velocidad de recesión $v$ frente a la distancia $d$ y, antes de ajustar, describe cualitativamente el patrón (\emph{¿qué ocurre con la velocidad cuando aumenta la distancia?}). Después realiza un ajuste lineal y obtiene la pendiente, que la web \textbf{no identifica inicialmente} como $H_0$: en su lugar se pregunta qué tipo de relación describen los datos, qué unidades tiene la pendiente, si todos los puntos caen sobre la recta, qué explica la dispersión y qué interpretación física cabría darle. La estrategia se inspira en la reconstrucción escolar del diagrama de Hubble \parencite{cardona2016}, y se modifica para asegurar que el patrón lo identifique el alumnado antes de introducir formalmente la ley. El producto es el diagrama de cada grupo, el valor experimental de su pendiente y su conjetura sobre el significado de esta.

\subsubsection{Fase 3: Explain --- construcción y formalización de la explicación}\label{sec:explain}

La fase \textit{Explain} comienza recuperando las explicaciones generadas en \textit{Explore}: cada grupo explica con sus palabras qué relación encontró entre distancia y velocidad y qué cree que representa la pendiente; las respuestas se comparan entre pares y se discuten. Solo después se introduce formalmente
\begin{equation}
v=H_0\,d,
\end{equation}
identificando la pendiente como la constante de Hubble y analizando sus unidades, $[H_0]=\mathrm{km\,s^{-1}\,Mpc^{-1}}$. La interpretación dimensional permite preguntar qué magnitud se obtiene al invertir $H_0$, de donde surge el \textbf{tiempo de Hubble},
\begin{equation}
t_H=\frac{1}{H_0},
\end{equation}
con la conversión explícita de unidades hasta años. Se insiste en que $t_H$ es una escala temporal y una estimación del orden de magnitud de la edad, pero no el cálculo cosmológico completo. La visualización de «rebobinado» de la expansión evita representar el Big Bang como una explosión en un punto privilegiado: muestra una distribución homogénea cuyas distancias comóviles disminuyen al retroceder en el tiempo, confrontando a la vez las ideas de «centro del universo» y de «Big Bang como explosión puntual». Se formaliza también la distinción entre movimiento peculiar y corrimiento al rojo cosmológico, aclarando que $v\simeq cz$ es una aproximación válida en el régimen de bajo $z$. El producto es una explicación escrita de qué representa $H_0$, por qué $1/H_0$ tiene unidades de tiempo y cuáles son los límites de esta estimación.

\subsubsection{Fase 4: Elaborate --- transferencia y análisis de los límites del modelo}\label{sec:elaborate}

La fase \textit{Elaborate} aplica lo aprendido a situaciones nuevas, sin convertirse en una exposición de temas nuevos: la materia oscura, la geometría global o los modelos alternativos quedan como ampliación opcional, fuera del núcleo evaluable. Se desarrolla en dos partes, igual que en la web.

\paragraph{El tamaño observable.} El alumnado calcula
\begin{equation}
R_H=\frac{c}{H_0},
\end{equation}
e identifica $R_H$ como el radio de Hubble, una escala característica de distancia ($\sim$14 mil millones de años-luz para $H_0\sim70$ km\,s$^{-1}$\,Mpc$^{-1}$), y la compara con el radio comóvil actual del universo observable ($\sim$46 mil millones de años-luz). La discrepancia se plantea como problema conceptual: si la luz ha viajado un tiempo comparable a la edad del universo, ¿cómo puede estar hoy su fuente a una distancia mucho mayor que $c\,t_0$? Una visualización interactiva muestra cualitativamente cómo el espacio se expande mientras la luz viaja; el objetivo no es resolver la integral cosmológica, sino comprender que $R_{\mathrm{observable}}\neq c\,t_0$ y que $c/H_0$ es una \textbf{escala}, no el radio exacto del universo observable.

\paragraph{Los límites de la medida.} La web proporciona dos valores plausibles de la constante de Hubble; el alumnado calcula el tiempo de Hubble correspondiente y razona qué ocurre con $1/H_0$ cuando aumenta $H_0$, qué valor produce una escala temporal mayor, qué significa que distintos métodos den valores algo diferentes y si ello implica que el conocimiento científico sea arbitrario. La actividad introduce cualitativamente la \textbf{tensión de Hubble} como ejemplo de incertidumbre, de contrastación entre métodos y del carácter revisable de la ciencia. Además, la estimación $t_H=1/H_0$ presupone una historia de expansión simplificada y la edad precisa requiere un modelo de la evolución de $H(t)$; la expansión acelerada y la energía oscura se introducen solo como contexto de esta limitación.

El producto de \textit{Elaborate} es la resolución argumentada de ambas partes y una breve reflexión sobre la diferencia entre una medida, una estimación y un modelo.

\subsubsection{Fase 5: Evaluate --- evaluación de la comprensión y reconstrucción del proceso}\label{sec:evaluate}

La fase \textit{Evaluate} no introduce contenidos nuevos. Comienza con el cuestionario conceptual final (\textit{postest}), equivalente al inicial y sin retroalimentación durante su aplicación, pues sus resultados sirven para comparar la comprensión antes y después. A continuación el alumnado elabora el \textbf{producto final}: \emph{explicar a una persona que no ha seguido la secuencia cómo sabemos que el universo se expande y cómo las observaciones permiten estimar su edad}. Para estructurarla se ofrece la cadena
\begin{equation}
\mathrm{espectro}\longrightarrow z\longrightarrow v\longrightarrow H_0\longrightarrow t_H,
\end{equation}
indicando en cada paso qué magnitud se observa directamente, cuál se calcula, cuál se infiere y qué aproximaciones se emplean. El producto final se evalúa con una rúbrica analítica que atiende a la corrección física, la coherencia del razonamiento, el reconocimiento de las aproximaciones y la claridad de la comunicación. El cuestionario \textit{pretest--postest} y las rúbricas se desarrollan en los instrumentos de evaluación.

\textbf{Flujo y cadena de preguntas (conexión entre fases).} La secuencia avanza como un \textbf{recorrido encadenado}: cada fase se enlaza con la siguiente mediante una \textbf{pregunta puente} (el nexo de unión), desde el problema inicial hasta la determinación de la edad y el tamaño.

\textbf{Plantilla común de las fichas de fase.} El desarrollo de cada fase se recoge en el Capítulo 4 (§4.2) mediante una ficha que comparte esta plantilla común:
\begin{itemize}
\item pregunta que la abre (de la cadena anterior);
\item objetivo(s) y OA vinculados (§3.4);
\item duración y agrupamiento;
\item qué sabemos e ideas previas;
\item contenido teórico de la fase;
\item desarrollo de la actividad (pasos);
\item interacción en la web;
\item idea errónea que confronta (§3.5);
\item espacio--tiempo que se trabaja (eje conceptual de §3.1);
\item materiales (§4.6);
\item producto de la fase;
\item indicadores de logro;
\item puente a la siguiente fase.
\end{itemize}
\textit{(La ficha de Evaluate añade, además, un apartado de diseño de evaluación y otro de proyección de la secuencia; véanse §4.2 y §4.3.)}

\textbf{Síntesis del recorrido.} La Figura~\ref{fig:flujo5e} resume las cinco fases como un recorrido encadenado, del problema inicial a la determinación de la edad y el tamaño, destacando la \emph{pregunta puente} (nexo) que enlaza cada fase con la siguiente.

\begin{figure}[htbp]
\centering
\begin{tikzpicture}[
  font=\footnotesize,
  node distance=2.6mm,
  fase/.style={rectangle, rounded corners, draw=naranja, line width=0.6pt, fill=naranja!10, text width=11.6cm, inner sep=1.8mm, align=left},
  extremo/.style={rectangle, rounded corners, draw=naranja, line width=0.9pt, fill=naranja, text=white, text width=11.6cm, inner sep=1.8mm, align=center},
  nexo/.style={rectangle, rounded corners, draw=black!25, dashed, fill=black!4, text width=9.6cm, inner sep=1.4mm, align=center, font=\scriptsize\itshape},
  flecha/.style={-{Stealth[length=2.4mm]}, draw=naranja, line width=0.9pt}
]
\node[extremo] (prob) {\textbf{Problema motriz:} ¿podemos determinar la edad y el tamaño del universo? ¿Cómo lo sabe la ciencia?};
\node[fase, below=of prob] (engage) {\textbf{\textcolor{naranja}{Engage}}: enmarcamos el reto, aclaramos qué son la \emph{edad} y el \emph{tamaño}, recorremos la historia de la cosmología y afloramos nuestras ideas; nos provoca la paradoja de la noche oscura (Olbers).};
\node[nexo, below=of engage] (n1) {nexo: si no podemos ir hasta allí, ¿qué información nos llega de las galaxias? $\rightarrow$ la luz};
\node[fase, below=of n1] (explore) {\textbf{\textcolor{naranja}{Explore}}: observamos el corrimiento al rojo, medimos velocidad y distancia y construimos el diagrama de Hubble $\rightarrow H_{0}$.};
\node[nexo, below=of explore] (n2) {nexo: hay un patrón ($v$ crece con $d$): ¿qué significa y qué número lo describe?};
\node[fase, below=of n2] (explain) {\textbf{\textcolor{naranja}{Explain}}: formalizamos $v=H_{0}d$; el corrimiento al rojo es estiramiento del espacio; invertimos $H_{0}\rightarrow$ edad ($t\approx 1/H_{0}$).};
\node[nexo, below=of explain] (n3) {nexo: ¿es exacta esa edad? ¿y hasta dónde podemos ver?};
\node[fase, below=of n3] (elaborate) {\textbf{\textcolor{naranja}{Elaborate}}: calculamos el tamaño observable y reconocemos sus \emph{límites} ---hasta dónde llegaremos a ver (el horizonte de sucesos)---; por qué $1/H_{0}$ es una aproximación; matiz de la energía oscura ($\Lambda$CDM y tensión de Hubble).};
\node[nexo, below=of elaborate] (n4) {nexo: ¿sabríamos explicar todo el proceso a otra persona?};
\node[fase, below=of n4] (evaluate) {\textbf{\textcolor{naranja}{Evaluate}}: demostramos que entendemos el \emph{proceso}, no solo el resultado: reconstruimos la cadena de la medida, leemos las ecuaciones y lo explicamos a otra persona (producto final), midiendo el avance con el cuestionario pre-post.};
\node[extremo, below=of evaluate] (res) {\textbf{Resultado:} la \textbf{edad} y el \textbf{tamaño} del universo observable, revisando en el camino las nociones de espacio y tiempo.};
\draw[flecha] (prob) -- (engage);
\draw[flecha] (engage) -- (n1);
\draw[flecha] (n1) -- (explore);
\draw[flecha] (explore) -- (n2);
\draw[flecha] (n2) -- (explain);
\draw[flecha] (explain) -- (n3);
\draw[flecha] (n3) -- (elaborate);
\draw[flecha] (elaborate) -- (n4);
\draw[flecha] (n4) -- (evaluate);
\draw[flecha] (evaluate) -- (res);
\end{tikzpicture}
\caption{Flujo de la secuencia 5E: del problema inicial a la determinación de la edad y el tamaño, con la \emph{pregunta puente} (nexo) que enlaza cada fase con la siguiente.}\label{fig:flujo5e}
\end{figure}

\subsection{Trazabilidad de las decisiones de diseño}

La Tabla~\ref{tab:trazabilidad5e} resume la procedencia de los principales elementos de la secuencia y el tipo de adaptación realizada.

\begin{table}[htbp]
\centering
\caption{Trazabilidad entre los antecedentes revisados y las actividades seleccionadas para la secuencia.}\label{tab:trazabilidad5e}
\begin{tabularx}{\textwidth}{|>{\raggedright\arraybackslash}p{2.4cm}|>{\raggedright\arraybackslash}X|>{\raggedright\arraybackslash}X|>{\raggedright\arraybackslash}X|}
\hline
\rowcolor{naranja}
{\color{white}\textbf{Recurso o antecedente}} & {\color{white}\textbf{Elemento considerado}} & {\color{white}\textbf{Decisión de diseño}} & {\color{white}\textbf{Adaptación realizada}} \\ \hline
Bybee et al. & Estructura del modelo 5E & Base organizativa de la secuencia & Se preserva la exploración antes de la formalización y la transferencia en \textit{Elaborate}. \\ \hline
NASA, \textit{Cosmic Times} & Historia de la cosmología, preguntas abiertas y actividades sobre la ley de Hubble & Uso parcial en \textit{Engage} y como antecedente de \textit{Explore} & Se reduce la carga histórica y se emplea sobre todo como elemento provocador y de naturaleza de la ciencia. \\ \hline
Las Cumbres Observatory, \textit{Measuring the Age of the Universe} & Espectros $\rightarrow$ redshift $\rightarrow$ velocidad $\rightarrow$ diagrama de Hubble $\rightarrow H_0 \rightarrow$ edad & Antecedente principal de la cadena observacional & Se elimina del núcleo la determinación de distancias con supernovas Ia, para centrar la actividad en el objetivo del TFM. \\ \hline
Cardona-Rodríguez et al. & Construcción de un diagrama de Hubble con datos astronómicos & Base metodológica para \textit{Explore} & Se incorpora el análisis en la propia plataforma y se pospone la formalización de la ley hasta \textit{Explain}. \\ \hline
NED-D / datos astronómicos públicos & Distancias independientes del corrimiento al rojo & Conjunto de datos reproducible & Se selecciona un subconjunto cerrado, documentado y almacenado localmente en la aplicación. \\ \hline
Literatura sobre concepciones en cosmología & Dificultades sobre expansión, Big Bang, redshift y escalas & Diseño de preguntas y conflictos cognitivos & Las concepciones se incorporan como distractores y como situaciones que deben explicarse durante las fases 5E. \\ \hline
\end{tabularx}
\end{table}
